% !TeX root = ../main.tex
% !TeX encoding = UTF-8

\chapter{The Governance Classification, Locating Your Organisation Across Five Dimensions}

Every organisation reading this book has a governance address. Most do not know what it is. A smaller number have a self-assessed view of it that, on examination, proves optimistic by two or three levels on at least one dimension. A very small number have had their address confirmed with forensic precision by someone whose professional interest is in accuracy rather than reassurance.

The Governance Classification is the diagnostic instrument this book has been building towards since Chapter 1. Everything in Parts One and Two established why the gap exists and what its consequences are. Chapters 11, 12, and 13 established what closing the gap requires operationally. This chapter establishes where a specific organisation sits in relation to that requirement, not in the abstract, but as a precise, five-dimension address that maps directly to the regulatory instruments in scope, the personal liabilities that are active, and the specific actions that would change the address.

The framework is derived from S.R. Ranganathan's Colon Classification theory, first formalised in 1933. Ranganathan recognised that complex subjects cannot be reduced to a single fixed hierarchy, knowledge must be expressed as a combination of independent dimensions, or facets, each with its own ordering logic. The combination produces an address, not a ranking. Applied to AI governance, this insight is structurally correct: an organisation can be at A4 and E1 simultaneously, named board-level accountability with no systematic AI inventory, and a composite score would collapse that into something that appears adequate. The governance address makes the gap structurally visible. A score conceals it.

\begin{quote}
\emph{A governance address is not a score. It is a location. An organisation at A4 and E1 simultaneously is not ``averaging'' to something adequate. It has named accountability with nothing to account for. The address reveals that gap. A score conceals it.}
\end{quote}

\scenebreak

\section{The Five Facets: What Each Dimension Measures}

The Governance Classification has five facets, each measured on a scale of one to five. The facets are independent, the level on any one facet does not constrain or determine the level on any other. The address is expressed as A·E·C·R·M, with the level on each facet stated in that order. An organisation at A3·E2·C3·R4·M2 has a specific, readable position on each dimension simultaneously.

\textbf{A, Accountability} \emph{Who owns AI decisions and can they be held to account?}

Accountability measures the presence and depth of named individual ownership of AI systems and decisions at board level. It is the facet most directly engaged by the FCA's SM\&CR Duty of Responsibility and the Companies Act's duty of care. A1 is no documented ownership at any level. A5 is the full Digital Alibi: all four elements present and active simultaneously at the output-moment level for every material AI-assisted decision.

A is the facet on which organisations most consistently overestimate their position. A board that has approved an AI strategy, named a Chief AI Officer, and received a quarterly AI governance report believes it is at A3 or A4. On examination, the accountability chain from AI system to board frequently lacks named individual attribution at the decision moment, it exists at the structural level but not at the decision level that SM\&CR requires. The gap between A3 (executive accountability exists; board ownership incomplete) and A4 (formal accountability chains documented at board level) is the gap the Vanishing Vendor scenario in Chapter 11 illustrates: structural accountability that has no internal anchor at the decision moment.

\textbf{E, Exposure} \emph{What AI risk is live, planned, or undetected?}

Exposure measures the completeness and currency of the AI systems inventory, including third-party and vendor-embedded AI. It is the facet most directly engaged by DORA's ICT risk management requirements and the EU AI Act's obligation to identify all high-risk AI systems in scope. E1 is no systematic inventory; shadow AI likely present. E5 is an exposure profile that updates continuously with emerging risks proactively identified.

E is the facet on which self-assessment is most unreliable. Chapter 12 established why: shadow AI and vendor-embedded AI are systematically absent from formal inventories. An organisation that has mapped its production AI systems but not its vendor-embedded AI, SaaS platform AI features, and internally adopted productivity AI is at E2, production systems mapped; third-party and shadow AI unmonitored, regardless of how comprehensive its formal inventory appears. The diagnostic correction most commonly required in advisory engagements is from E3 to E2 when vendor-embedded AI is properly surfaced.

\textbf{C, Control} \emph{What governance structures and oversight exist?}

Control measures the operational state and independence of AI governance structures. It distinguishes between governance that exists on paper (C1), governance that is operational but not independently tested (C3), and governance that is independently verified with evidence retained for regulators (C4). C5 is governance that evolves with the AI landscape, with the board receiving regular assurance as a standard operational discipline.

C is the facet most directly engaged by ISO 42001 certification, the EU AI Act's Article 17 quality management requirements, and Lloyd's of London's syndicate expectations. The most common gap in the C facet is between C3 (governance framework operational; not yet independently tested) and C4 (controls independently verified; evidence retained). An organisation that has built a governance architecture but has not subjected it to independent testing, including the retrievability test from Chapter 13, is at C3. The chapter's four-question retrievability test is, in effect, a C4 readiness diagnostic.

\textbf{R, Regulation} \emph{Which obligations apply and are they met?}

Regulation measures the completeness of the organisation's identification of, and compliance with, the AI governance obligations that apply to it. R1 is EU AI Act, FCA, and ICO obligations not identified or assessed. R5 is ahead of current obligations, positioned for emerging requirements before they are mandated.

R is the facet most directly affected by the 2 August 2026 EU AI Act deadline for high-risk AI system operators. An organisation that has not completed its EU AI Act gap analysis, identifying which of its AI systems fall within the high-risk classification and what obligations apply to each, is at R2 or below. An organisation with an active compliance programme that has not yet achieved audit-ready status is at R3. The 2 August 2026 deadline set the R4 threshold: obligations materially met, documented, and defensible to regulators.

\textbf{M, Maturity} \emph{How far has governance developed and in which direction?}

Maturity measures the depth and direction of governance development, not as an absolute state but as a trajectory. M1 is a board with no structured view of AI governance requirements. M5 is an organisation where governance capability is a competitive and reputational asset, embedded as an operational discipline rather than a compliance programme.

M is the facet that measures whether governance is improving or static. An organisation at M3 (governance work underway; not yet embedded organisationally) that is actively progressing towards M4 is in a different position from one at M3 that has been at M3 for two years without material progress. The Governance Classification does not capture trajectory, it captures position. But the gap assessment that follows from the classification reveals the distance to the required address, which is the measure of urgency.

\scenebreak

\section{What Each Level Means in Regulatory and Legal Terms}

The following table maps each level across all five facets to its regulatory position, what instruments are in scope, what exposure is active, and what a regulator or court would see at that address.

\begin{longtable}[]{@{}
  >{\raggedright\arraybackslash}p{(\columnwidth - 4\tabcolsep) * \real{0.1401}}
  >{\raggedright\arraybackslash}p{(\columnwidth - 4\tabcolsep) * \real{0.0586}}
  >{\raggedright\arraybackslash}p{(\columnwidth - 4\tabcolsep) * \real{0.8013}}@{}}
\toprule\noalign{}
\begin{minipage}[b]{\linewidth}\raggedright
\textbf{Level}
\end{minipage} & \begin{minipage}[b]{\linewidth}\raggedright
\textbf{Address}
\end{minipage} & \begin{minipage}[b]{\linewidth}\raggedright
\textbf{Regulatory position}
\end{minipage} \\
\midrule\noalign{}
\endhead
\bottomrule\noalign{}
\endlastfoot
\textbf{Level 1 Unstructured Exposure} & \textbf{A1·E1·C1·R1·M1} & No formal compliance standard in scope. General company law and common law duty of care only. Neither was written with AI in mind. Absence of governance is itself a finding under FCA SM\&CR and EU AI Act Article 14. \\
\textbf{Level 2 Voluntary Principles Only} & \textbf{A2·E2·C2·R2·M2} & OECD AI Principles, UNESCO Recommendation, UK AI White Paper, EU AI Pact. All voluntary. No enforcement weight. No legal protection if an incident occurs. Intent documented; no evidence. \\
\textbf{Level 3 Entry-Level Compliance} & \textbf{A3·E3·C3·R3·M3} & First level at which formal compliance obligations are engaged: GDPR Article 22, ICO AI Guidance, EU AI Act GPAI, FCA SM\&CR minimum. Programme in place; compliance active but not yet independently evidenced. \\
\textbf{Level 4 Formally Verified} & \textbf{A4·E4·C4·R4·M4} & ISO 42001:2023 and ISO 27001:2022 certifiable. EU AI Act High Risk compliant (2 August 2026 deadline). DORA and NIS2 obligations met. Defensible under regulatory scrutiny; evidence retained and retrievable. \\
\textbf{Level 5 Governance as Competitive Asset} & \textbf{A5·E5·C5·R5·M5} & All current obligations met and evidenced. Full Digital Alibi active across all material decisions. Ahead of EU AI Act revisions and AI Liability Directive. Governance is a demonstrable asset in PE due diligence, procurement, and board assurance. \\
\end{longtable}

Two observations about this table are operationally critical. First, Level 2 provides no legal protection. An organisation that has adopted the OECD AI Principles, published an AI ethics statement, and joined the EU AI Pact is at approximately A2·E2·C2·R2·M2. That position has intent and no evidence. When an incident occurs, the voluntary commitments will be adduced as evidence that the organisation knew the governance standard it had not met, not as evidence that it had met it. Second, Level 3 is the convergence point, the minimum at which any formal compliance obligation is engaged. The large majority of formal compliance standards converge at A3·E3·C3·R3·M3 as their baseline. Below this address, an organisation has not met any formal standard.

\begin{quote}
\emph{Level 2 is not a governance position. It is a liability position with a principles overlay. When an incident occurs, the voluntary commitments will be adduced as evidence that the organisation knew the standard it had not met. Voluntary principles without operational evidence are worse than silence.}
\end{quote}

\scenebreak

\section{The A3·E3·C3·R3·M3 Convergence Point}

The most significant structural finding in the AI governance standards landscape is the convergence of formal compliance standards at A3·E3·C3·R3·M3. GDPR Article 22, ICO AI Guidance, EU AI Act General Purpose AI obligations, FCA SM\&CR baseline, and the UK AI Safety Framework all require this address or its constituent elements. This convergence is not coincidental. It represents the structural minimum at which governance has moved from intention to operational reality.

At A3, executive accountability exists and board ownership is incomplete. At E3, all AI systems have been identified but risk has not been formally assessed. At C3, a governance framework is operational but not independently tested. At R3, a compliance programme is active with gaps remaining but not yet audit-ready. At M3, governance work is underway but not yet embedded. This is the address at which an organisation has made a genuine governance commitment and has not yet completed the work that commitment requires.

Below A3·E3·C3·R3·M3 on any facet, an organisation is not in the scope of any formal compliance standard. Above it, the standards begin to diverge, some requiring higher levels on specific facets, others satisfied by the convergence point. The gap between where most organisations self-assess and where they actually sit is typically at or below the convergence point on the E facet, because shadow AI and vendor-embedded AI are not captured in the formal inventory that produces the self-assessment.

\scenebreak

\section{Sector Ceilings, Where Regulated Organisations Must Reach}

For organisations in regulated sectors, the convergence point is the floor, not the ceiling. The combined regulatory stack applicable to financial services, insurance, aviation, healthcare, and the public sector drives the governance address requirement above A3·E3·C3·R3·M3 on multiple facets simultaneously. The sector ceiling is the governance destination: not where the organisation is, but where the regulatory environment and fiduciary obligations require it to be.

\begin{longtable}[]{@{}
  >{\raggedright\arraybackslash}p{(\columnwidth - 6\tabcolsep) * \real{0.1345}}
  >{\raggedright\arraybackslash}p{(\columnwidth - 6\tabcolsep) * \real{0.1170}}
  >{\raggedright\arraybackslash}p{(\columnwidth - 6\tabcolsep) * \real{0.1111}}
  >{\raggedright\arraybackslash}p{(\columnwidth - 6\tabcolsep) * \real{0.6374}}@{}}
\toprule\noalign{}
\begin{minipage}[b]{\linewidth}\raggedright
\textbf{Sector}
\end{minipage} & \begin{minipage}[b]{\linewidth}\raggedright
\textbf{Required address}
\end{minipage} & \begin{minipage}[b]{\linewidth}\raggedright
\textbf{Pressure facets}
\end{minipage} & \begin{minipage}[b]{\linewidth}\raggedright
\textbf{Primary driver}
\end{minipage} \\
\midrule\noalign{}
\endhead
\bottomrule\noalign{}
\endlastfoot
\textbf{Financial Services} & \textbf{A4·E4·C4·R4·M4} & A + E & FCA SM\&CR personal accountability (A4) + DORA third-party AI inventory (E4) \\
\textbf{Insurance (Lloyd's)} & \textbf{A4·E4·C4·R4·M4} & A + C & SM\&CR board accountability chain (A4) + independently verified controls for Lloyd's and EU AI Act Art.17 (C4) \\
\textbf{Aviation} & \textbf{A4·E4·C4·R4·M4} & C + R & Operational control integrity (C4) + CAA/EU AI Act/NIS2 multi-framework compliance (R4) \\
\textbf{Healthcare} & \textbf{A4·E3·C4·R4·M3} & A + C & Named clinical accountability for CQC and EU AI Act High Risk (A4) + evidence-based controls (C4) \\
\textbf{Public Sector} & \textbf{A4·E3·C3·R3·M3} & A & UK AI Safety Framework contestability requirement drives A4; other facets at 3 \\
\end{longtable}

The pressure facets, the dimensions on which each sector's regulatory stack places the most demanding requirements, are the governance investment priorities for any organisation in that sector. A financial services organisation must reach A4 and E4 before the other facets, because those are the dimensions on which SM\&CR and DORA create the most direct personal and institutional liability. A healthcare organisation must reach A4 and C4 before the others, because clinical accountability and evidence-based controls are the dimensions on which CQC and EU AI Act High Risk create the most direct patient safety and regulatory exposure.

The gap between a confirmed starting address and the sector ceiling is the governance investment required. An organisation in financial services that confirms its address at A2·E2·C2·R2·M2 is not simply ``at an early stage.'' It is absent from the regulatory landscape on all five dimensions in a sector where the combined requirement is A4·E4·C4·R4·M4. That is a specific, named, and quantified liability, not a maturity journey.

\scenebreak

\section{Using the Governance Classification as a Diagnostic Instrument}

The Governance Classification produces an address, not an assessment. The address is neutral, it describes where an organisation is, not whether that position is adequate. The adequacy question is answered by comparing the address to the sector ceiling and the regulatory obligations in scope. The gap between the two is the governance exposure.

The diagnostic process has three stages. The first is self-assessment: using the level descriptors for each facet to form an initial view of the organisation's position. Self-assessment is always provisional. The primary value of self-assessment is to establish the starting hypothesis that the diagnostic will confirm or correct.

The second stage is confirmation: subjecting the self-assessed address to forensic examination against the evidence that each level descriptor requires. An organisation that self-assesses at A3 must confirm that named executive accountability exists and is documented at the level the descriptor requires, not merely that a governance structure is in place. An organisation that self-assesses at E3 must confirm that all AI systems, including vendor-embedded AI and shadow AI, have been identified, not merely that a formal AI inventory exists. The gap between self-assessed and confirmed address is itself a governance finding. It tells the board what it believed about its governance position and what the evidence actually shows.

The third stage is gap analysis: mapping the confirmed address against the sector ceiling and the specific regulatory obligations in scope, identifying the facets on which the gap is widest, and prioritising the governance investments required to close it. The gap analysis is not a maturity roadmap. It is a liability map, showing where the organisation is most exposed, in regulatory and legal terms, at its confirmed address.

\begin{longtable}[]{@{}
  >{\raggedright\arraybackslash}p{(\columnwidth - 0\tabcolsep) * \real{1.0000}}@{}}
\toprule\noalign{}
\begin{minipage}[b]{\linewidth}\raggedright
\textbf{READING A GOVERNANCE ADDRESS, A WORKED EXAMPLE}
\end{minipage} \\
\midrule\noalign{}
\endhead
\bottomrule\noalign{}
\endlastfoot
Confirmed address: A2·E2·C3·R3·M2. Sector: financial services. Required ceiling: A4·E4·C4·R4·M4. \\
A2: Technical teams own AI systems; board accountability is absent. SM\&CR Duty of Responsibility is not met. Personal liability exposure for the gap between the AI decision and the named board-level accountable individual is uninsured. \\
E2: Production systems mapped; third-party and shadow AI unmonitored. DORA ICT third-party risk requirements are not met. The AI inventory that accountability chains must attach to is incomplete. \\
C3: Governance framework operational; not yet independently tested. The retrievability test from Chapter 13 has not been passed. EU AI Act Article 17 quality management requirements are not met. \\
R3: Compliance programme active; gaps remain; not audit-ready. The EU AI Act obligations in force since 2 August 2026 are not met at R3. Personal and institutional penalties of up to €15 million or 3\% of global annual turnover remain as active exposure. \\
M2: Governance need understood at executive level; no structured programme. The board has discussed AI governance but has not built the architecture to deliver it. \\
Summary: this organisation has two priority investments, A and E, before any other facet. Without A4, the SM\&CR liability is personal and active. Without E4, there is no complete inventory for the accountability chain to attach to. C, R, and M cannot be adequately addressed until A and E are confirmed. The investment sequence is determined by the regulatory stack, not by the board's preference for which dimension to tackle first. \\
\end{longtable}

\scenebreak

\section{A5: The Digital Alibi as a Governance Address}

The A facet's fifth level, A5, is the governance address that encodes the full Digital Alibi standard. It is the only level in the Governance Classification whose definition is derived from the four Digital Alibi elements rather than from a regulatory instrument or certification standard.

A5 requires four conditions to be simultaneously present and active: information provenance, identifiable, versioned sources for each material AI-assisted decision, retrievable on demand; a retrieval log, a documented record of what the system searched for and what it did not find, tied to the decision timestamp; a human observability record, a documented record of who saw which outputs at which stage of the decision process and what action they took; and a retrievability guarantee, all of the above is not merely archived but actively retrievable within a timeframe that would satisfy a regulator or court.

A5 is not achievable through any single governance framework or certification. ISO 42001 reaches A4 on the accountability dimension if fully implemented. The EU AI Act's High Risk obligations, fully met, reach A4. No existing regulatory instrument or certification standard requires A5 because no existing instrument requires the full Digital Alibi at the output-moment level for every material decision. A5 is the governance address that this book has established as the defensible standard, not the regulatory minimum, but the evidentiary standard that a challenge from a regulator, litigant, or shareholder will require.

\begin{quote}
\emph{A5 is not the regulatory minimum. It is the evidentiary standard that a challenge will require. The regulatory minimum (A4) documents that governance exists. A5 demonstrates that governance operated, at this decision, at this moment, with this information picture. The gap between A4 and A5 is the gap between documentation and defensibility.}
\end{quote}

No single governance entity currently reaches A5 alone. The Digital Alibi architecture requires the combination of governance protocol at the board level and a decision-level record architecture at the output-moment level. Chapter 20 describes the Alibi Record, the technical and governance implementation of A5, as the specific deliverable of the Protocol that Part Four introduces.

\scenebreak

\section{The Cardinal Rule: Score Only What Is Evidenced}

The Governance Classification is a diagnostic instrument, not a self-certification. Its value depends entirely on the integrity of the assessment. An organisation that scores itself at A4 because it has an SM\&CR accountability map, without confirming that the map is current at the output-moment level, that named individuals are documented contemporaneously for each material AI system, and that the chain survives forensic review, has produced a self-assessed address that will not correspond to the confirmed address.

The cardinal rule of the Governance Classification is: score only what is evidenced. If a facet level requires named individual accountability at the output-moment level and the evidence is a structural accountability map reviewed annually, the evidenced level is A3, not A4. If a facet level requires all AI systems to be identified including vendor-embedded AI, and the evidence is a formal AI inventory that was not built to include SaaS platform AI features, the evidenced level is E2, not E3. The gap between what is claimed and what is evidenced is not a governance failure, it is the finding that the diagnostic exists to produce.

The most valuable output of the Governance Classification is not a high address. It is an accurate address. A board that knows its confirmed address at A2·E2·C3·R3·M2 knows precisely where its liability is concentrated, what the gap to its sector ceiling is, and what investments would change its address. A board that self-assesses at A3·E3·C3·R3·M3 without confirming those levels against evidence has a confidence it has not earned and a liability it has not measured.

\scenebreak

\section{The Problem That Gets Exponentially Worse: Agentic AI}

Chapter 14 has given the board its diagnostic instrument. The Governance Classification locates an organisation in the AI governance landscape with precision, confirming where it is, what the gap to its sector ceiling is, and what the priority investments are. For the AI systems most organisations currently deploy, that diagnostic is sufficient to scope the governance work required.

Chapter 15 examines the form of AI deployment that the Governance Classification cannot fully address with the tools described so far, agentic AI systems that operate with significant autonomy across extended decision sequences, taking actions rather than producing recommendations, and creating accountability chains whose links are generated by the system rather than by human governance decisions. Agentic AI does not change the Governance Classification framework. It makes the gap at every facet structurally harder to close, and it does so in ways that most boards have not yet encountered in their AI governance work.
